% GENERATED by scripts/registry/gen_registry.py. Do not edit.
% Inputs (sha256):
%   registry/notation.json b47459857bd3809f8fa1776487a8350b2b18939e011c1610894f704d04adf4c2
%   registry/definitions.json d636fa56008403ceb21229ab37e47df58a10d3661453da7f1e4176955e8e1f09
%   registry/conventions.json f5380d7a56bf2f81fcc1f5f1d6f4ddd2cbee86ee1aa8fd61e1562860e592969c
%   registry/translation.json 5aff431c79e6c43413aae14995775c5eca8b7df42c6134c2cdb832d7543b63b1
% One declaration per definition, in the interface of latex/tpc-course.sty:
%   \tpcdeclaredef{ID}{depth}{Titel}{Title}{deutscher Text}{English text}
% tpc-course.sty turns \tpcdef{D17} into a numbered definition in the current
% language and checks the depth. The \providecommand lines below are a plain
% fallback so that this file also works without tpc-course.sty: there
% \tpcdef{D17} prints the name in bold and the body, in the language chosen by
% \tpcifenglish (German when that macro is absent). An unknown ID is an error.
\makeatletter
\providecommand{\tpcifenglish}[2]{#2}
\providecommand{\tpcdeclaredef}[6]{%
  \@namedef{tpcdef@depth@#1}{#2}%
  \long\@namedef{tpcdef@title@#1}{\tpcifenglish{#4}{#3}}%
  \long\@namedef{tpcdef@body@#1}{\tpcifenglish{#6}{#5}}}
\providecommand{\tpcdef}[1]{%
  \@ifundefined{tpcdef@body@#1}%
    {\PackageError{tpc-definitions}{Definition `#1' is not in the registry}{See registry/definitions.json.}}%
    {\par\noindent\textbf{Definition #1 (\@nameuse{tpcdef@title@#1}).} \@nameuse{tpcdef@body@#1}\par}}
\makeatother
% D1 [blessed-provisional]
\tpcdeclaredef{D1}{kern}{Testdefinition}{test definition}%
  {Ein Zustandsraum $\StSp$.}%
  {Ein Zustandsraum $\StSp$.}
% D2 [blessed-provisional]
\tpcdeclaredef{D2}{fundament}{Testdefinition}{test definition}%
  {Eine Ordnungseinheit $\OrdUnit$.}%
  {Eine Ordnungseinheit $\OrdUnit$.}
\endinput
